(1) 两体库仑问题将相对坐标运动简化成质量 \(\mu=m_eM/(m_e+M)\) 的有效单体问题，质心自由运动可与内部束缚分离。有限核质量使 \(\mu<m_e\)，这是同位素谱线移动的最低阶非相对论来源。
(2) 库仑势 \(V(r)=-Ze^2/(4\pi\varepsilon_0r)\)，能级 \(E_n=-\mu Z^2e^4/[2(4\pi\varepsilon_0)^2\hbar^2n^2]\)，\(n=1,2,\ldots\)。相应尺度 \(a_Z=4\pi\varepsilon_0\hbar^2/(\mu Ze^2)\)。质量增大会增大束缚能绝对值、减小轨道尺度；写成 \(m_e\) 的表达式需在无限核质量近似才成立。
(3) 能量守恒 \(h\nu=E_{n_i}-E_{n_f}\) 的符号取发光光子正值，故 \(\nu=[\mu Z^2e^4/(2h(4\pi\varepsilon_0)^2\hbar^2)](1/n_f^2-1/n_i^2)\)。当 \(M\to\infty\) 时 \(\mu\to m_e\)，恢复通常的氢样 Rydberg 频率表达式。
(4) 在 \(M\gg m_e\) 时 \(\mu=m_e(1-m_e/M+O(m_e^2/M^2))\)。相同跃迁能级之间频率正比 \(\mu\)，所以 \([\nu(M_2)-\nu(M_1)]/\nu(\infty)\simeq m_e(1/M_1-1/M_2)\)。若 \(M_2>M_1\)，频率升高而波长减小；这里只含约化质量项，不包含核体积、精细结构和超精细结构。
